% =====================================================================
%  courses/aritmetica/semana_01_01_leyes_logica_propocicional/teoria.tex — teoría: leyes de la lógica proposicional y circuitos lógicos
% ---------------------------------------------------------------------
%  Curso: Aritmética | Semana: 01 | Unidad 1: Fundamentos Numéricos
% =====================================================================

% --- PORTADA DE CLASE ------------------------------------------------------
\portadaclase{Leyes de la L\'{o}gica Proposicional}{01}{1}{Fundamentos de L\'{o}gica}

% --- MOTIVACIÓN ------------------------------------------------------------
\section*{¿Por qué estudiar este tema?}
\begin{tcolorbox}[
  enhanced, colback=AzulClaro!30, colframe=AzulMedio,
  arc=3pt, boxrule=0.6pt, left=10pt, right=10pt, top=6pt, bottom=6pt
]
\textit{Las leyes de la lógica proposicional son la base del razonamiento
formal en matemáticas y están presentes en exámenes de admisión de la UNI,
San Marcos, PUCP y academias Aduni, Lumbreras y Trilce. Dominar las
equivalencias lógicas permite simplificar expresiones complejas en segundos,
una habilidad clave para resolución rápida en preselección.}
\end{tcolorbox}
\vspace{0.5cm}

% --- SECCIÓN 1: PROPOSICIONES EQUIVALENTES ---------------------------------
\section{Proposiciones Equivalentes}

\subsection{Definición}

\begin{cajadef}
  \textbf{Proposiciones equivalentes:} Dos proposiciones $A$ y $B$ son
  equivalentes cuando la bicondicional $A \leftrightarrow B$ es una
  \textbf{tautología}. Se denota:
  \[
    A \equiv B \qquad \text{(``$A$ es equivalente a $B$'')}
  \]
  Es decir, $A$ y $B$ tienen \textit{exactamente} los mismos valores de
  verdad en toda circunstancia posible.
\end{cajadef}

\begin{cajatip}
  Una equivalencia lógica \textbf{no} es simplemente que ambas sean
  verdaderas en un caso; deben coincidir en \textit{todas} las filas de
  su tabla de verdad. Verificar con la bicondicional: si siempre es
  verdadera $(\mathbf{V})$, entonces son equivalentes.
\end{cajatip}

% --- SECCIÓN 2: LEYES DEL ÁLGEBRA DE PROPOSICIONES -------------------------
\section{Leyes del \'{A}lgebra de Proposiciones}

A continuación se presentan las 11 leyes fundamentales con sus expresiones
formales. El dominio de estas leyes permite \textbf{simplificar} esquemas
moleculares complejos en pasos sistemáticos.

\subsection{Ley 1 — Doble Negación (Involutiva)}

\begin{cajaprop}[1 — Doble Negaci\'{o}n]
  Negar dos veces una proposición devuelve la proposición original:
  \[
    \sim(\sim p) \equiv p
  \]
\end{cajaprop}

\subsection{Ley 2 — Idempotencia}

\begin{cajaprop}[2 — Idempotencia]
  Operar una proposición consigo misma no cambia su valor:
  \[
    p \wedge p \equiv p \qquad\qquad p \vee p \equiv p
  \]
\end{cajaprop}

\subsection{Ley 3 — Conmutativa}

\begin{cajaprop}[3 — Conmutativa]
  El orden de las proposiciones en la conjunción o disyunción no altera
  el resultado:
  \[
    p \wedge q \equiv q \wedge p \qquad\qquad p \vee q \equiv q \vee p
  \]
\end{cajaprop}

\subsection{Ley 4 — Asociativa}

\begin{cajaprop}[4 — Asociativa]
  El agrupamiento de proposiciones en conjunciones o disyunciones puede
  modificarse libremente:
  \[
    p \wedge (q \wedge r) \equiv (p \wedge q) \wedge r
  \]
  \[
    p \vee (q \vee r) \equiv (p \vee q) \vee r
  \]
\end{cajaprop}

\subsection{Ley 5 — Distributiva}

\begin{cajaprop}[5 — Distributiva]
  La conjunción distribuye sobre la disyunción y viceversa:
  \[
    p \wedge (q \vee r) \equiv (p \wedge q) \vee (p \wedge r)
  \]
  \[
    p \vee (q \wedge r) \equiv (p \vee q) \wedge (p \vee r)
  \]
\end{cajaprop}

\subsection{Ley 6 — De Morgan}

\begin{cajaprop}[6 — De Morgan]
  La negación de una conjunción equivale a la disyunción de las
  negaciones, y viceversa:
  \[
    \sim(p \wedge q) \equiv \sim p \vee \sim q
  \]
  \[
    \sim(p \vee q) \equiv \sim p \wedge \sim q
  \]
\end{cajaprop}

\begin{cajaadvert}
  \begin{itemize}
    \item \textbf{Error frecuente:} Al aplicar De Morgan, muchos estudiantes
      olvidan cambiar el conector: $\sim(p \wedge q) \neq \sim p \wedge \sim q$.
      \textit{El $\wedge$ debe convertirse en $\vee$.}
    \item \textbf{Regla mnemotécnica:} ``Entro con negación, cambio el
      conector y niego cada proposición.''
  \end{itemize}
\end{cajaadvert}

\subsection{Ley 7 — De la Condicional}

\begin{cajaprop}[7 — De la Condicional]
  La condicional puede reescribirse en términos de disyunción, y también
  posee una forma transpuesta equivalente:
  \[
    p \rightarrow q \equiv \sim p \vee q
  \]
  \[
    p \rightarrow q \equiv \sim q \rightarrow \sim p \qquad
    \text{(Transposici\'{o}n o Contrarrecíproca)}
  \]
\end{cajaprop}

\subsection{Ley 8 — De la Bicondicional}

\begin{cajaprop}[8 — De la Bicondicional]
  La bicondicional equivale a la conjunción de dos condicionales en
  sentidos opuestos:
  \[
    p \leftrightarrow q \equiv (p \rightarrow q) \wedge (q \rightarrow p)
  \]
\end{cajaprop}

\subsection{Ley 9 — Absorción}

\begin{cajaprop}[9 — Absorci\'{o}n]
  Una proposición ``absorbe'' combinaciones que la contienen implícitamente:
  \begin{align*}
    p \wedge (p \vee q) &\equiv p \\[4pt]
    p \vee (p \wedge q) &\equiv p \\[4pt]
    p \wedge (\sim p \vee q) &\equiv p \wedge q \\[4pt]
    p \vee (\sim p \wedge q) &\equiv p \vee q
  \end{align*}
\end{cajaprop}

\begin{cajatip}
  La ley de absorción es la más utilizada en simplificación de circuitos
  lógicos y reducción de esquemas moleculares. Identifique siempre primero
  si una subexpresión contiene la proposición principal.
\end{cajatip}

\subsection{Ley 10 — Complemento}

\begin{cajaprop}[10 — Complemento]
  La disyunción de una proposición con su negación es siempre verdadera
  (tautología); su conjunción es siempre falsa (contradicción):
  \[
    p \vee \sim p \equiv \mathbf{V} \qquad\qquad p \wedge \sim p \equiv \mathbf{F}
  \]
\end{cajaprop}

\subsection{Ley 11 — Identidad}

\begin{cajaprop}[11 — Identidad]
  Las constantes $\mathbf{V}$ (verdad) y $\mathbf{F}$ (falsedad) actúan
  como elementos neutros y absorbentes:
  \[
    P \vee \mathbf{V} \equiv \mathbf{V} \qquad\quad
    P \vee \mathbf{F} \equiv P
  \]
  \[
    P \wedge \mathbf{V} \equiv P \qquad\quad
    P \wedge \mathbf{F} \equiv \mathbf{F}
  \]
\end{cajaprop}

% --- SECCIÓN 3: CIRCUITOS LÓGICOS ------------------------------------------
\section{Circuitos L\'{o}gicos}

\subsection{Definición}

\begin{cajadef}
  Un \textbf{circuito lógico} es un arreglo de interruptores denominados
  \textbf{compuertas lógicas}, donde cada compuerta tiene asociado un valor
  de verdad. El estado del circuito (abierto o cerrado) determina si la
  corriente pasa o no, lo cual modela el valor de verdad de la proposición
  compuesta correspondiente.
\end{cajadef}

\subsection{Tipos de circuito}

\begin{cajaprop}[Circuito en Serie — Conjunci\'{o}n ($\wedge$)]
  Los interruptores $p$ y $q$ están en serie: la corriente pasa
  \textbf{solo si ambos} están cerrados (ambos verdaderos).
  \[
    \text{Circuito en serie} \equiv p \wedge q
  \]
  \begin{center}
  \begin{tikzpicture}[
    switch/.style={
      draw=AzulPizarra, fill=AzulClaro!40,
      minimum width=1.2cm, minimum height=0.5cm,
      rounded corners=2pt, font=\sffamily\small
    },
    wire/.style={draw=AzulPizarra, thick},
    node distance=0.5cm
  ]
    \coordinate (L) at (0,0);
    \coordinate (R) at (6.5,0);
    \draw[wire] (L) -- ++(0.6,0);
    \node[switch] (P) at (1.8,0) {$p$};
    \draw[wire] (P.east) -- ++(0.5,0);
    \node[switch] (Q) at (3.8,0) {$q$};
    \draw[wire] (Q.east) -- ++(0.6,0);
    \draw[wire] (L) -- ++(0,0);
    \node[right, font=\sffamily\small\color{AzulPizarra}]
      at (5.8,0) {$p \wedge q$};
    \filldraw[AzulPizarra] (L) circle (2pt);
    \filldraw[AzulPizarra] (5.5,0) circle (2pt);
  \end{tikzpicture}
  \end{center}
\end{cajaprop}

\begin{cajaprop}[Circuito en Paralelo — Disyunci\'{o}n ($\vee$)]
  Los interruptores $p$ y $q$ están en paralelo: la corriente pasa
  \textbf{si al menos uno} está cerrado (al menos uno verdadero).
  \[
    \text{Circuito en paralelo} \equiv p \vee q
  \]
  \begin{center}
  \begin{tikzpicture}[
    switch/.style={
      draw=AzulPizarra, fill=AzulClaro!40,
      minimum width=1.2cm, minimum height=0.5cm,
      rounded corners=2pt, font=\sffamily\small
    },
    wire/.style={draw=AzulPizarra, thick}
  ]
    \coordinate (A) at (0,0);
    \coordinate (B) at (5.0,0);
    % Nodos de derivación
    \coordinate (LA) at (0.8, 0);
    \coordinate (RA) at (4.2, 0);
    % Rama superior
    \draw[wire] (LA) -- ++(0, 0.7) -- ++(0.6,0);
    \node[switch] (P) at (2.3, 0.7) {$p$};
    \draw[wire] (P.east) -- ++(0.6,0) -- ++(0,-0.7);
    % Rama inferior
    \draw[wire] (LA) -- ++(0,-0.7) -- ++(0.6,0);
    \node[switch] (Q) at (2.3,-0.7) {$q$};
    \draw[wire] (Q.east) -- ++(0.6,0) -- ++(0, 0.7);
    % Línea principal
    \draw[wire] (A) -- (LA);
    \draw[wire] (RA) -- (B);
    \filldraw[AzulPizarra] (A)  circle (2pt);
    \filldraw[AzulPizarra] (B)  circle (2pt);
    \filldraw[AzulPizarra] (LA) circle (2pt);
    \filldraw[AzulPizarra] (RA) circle (2pt);
    \node[right, font=\sffamily\small\color{AzulPizarra}]
      at (5.1,0) {$p \vee q$};
  \end{tikzpicture}
  \end{center}
\end{cajaprop}

\subsection{Lectura de circuitos mixtos}

\begin{cajaejemplo}[1 — Lectura de esquema molecular desde un circuito]
  \textbf{Enunciado:} Determinar el esquema molecular para el circuito
  con dos ramas en paralelo, donde la rama superior tiene $p$ y $q$ en
  serie, y la rama inferior tiene $\sim p$ y $t$ en serie; todo ello
  compuesto con una sección adicional en serie con $r$.

  \noindent\textbf{Análisis:}
  \begin{itemize}
    \item Rama superior paralela: $p \wedge q$
    \item Rama inferior paralela: $\sim p \wedge t$
    \item Bloque paralelo completo: $(p \wedge q) \vee (\sim p \wedge t)$
    \item Sección en serie con $r$ al final:
      $\bigl[(p \wedge q) \vee (\sim p \wedge t)\bigr] \wedge r$
  \end{itemize}
  \begin{tcolorbox}[colback=AzulClaro!50, colframe=AzulPizarra,
    arc=2pt, boxrule=0.7pt]
    \[
      \boxed{\bigl[(p \wedge q) \vee (\sim p \wedge t)\bigr] \wedge r}
    \]
  \end{tcolorbox}
\end{cajaejemplo}

% --- SECCIÓN 4: ESTRATEGIA DE SIMPLIFICACIÓN -------------------------------
\section{Estrategia de Simplificaci\'{o}n Paso a Paso}

\begin{cajaformula}
  \renewcommand{\arraystretch}{1.5}
  \begin{center}
  \begin{tabular}{@{}cl@{}}
    \toprule
    \rowcolor{AzulClaro!60}
    \textbf{Paso} & \textbf{Acción recomendada} \\
    \midrule
    1 & Eliminar condicionales con Ley~7:
        $p \rightarrow q \equiv \sim p \vee q$ \\
    2 & Eliminar bicondicionales con Ley~8 \\
    3 & Ingresar negaciones con De Morgan (Ley~6) \\
    4 & Aplicar doble negación (Ley~1) para limpiar \\
    5 & Aplicar distributiva (Ley~5) si es necesario \\
    6 & Identificar y aplicar absorción (Ley~9) \\
    7 & Reducir con complemento (Ley~10) e identidad (Ley~11) \\
    \bottomrule
  \end{tabular}
  \end{center}
\end{cajaformula}

% --- EJEMPLOS RESUELTOS EN TEORÍA ------------------------------------------
\section{Ejemplos resueltos}

\begin{cajaejemplo}[2 — Nivel básico: Simplificación con De Morgan y absorción]
  \textbf{Enunciado:} Simplificar la expresión:
  \[
    (p \vee q) \vee \bigl[(\sim p \wedge \sim q) \vee p\bigr]
  \]
  \noindent\textbf{Método:} Leyes de absorción, complemento e identidad.
  \medskip

  \begin{align*}
    &(p \vee q) \vee \bigl[(\sim p \wedge \sim q) \vee p\bigr] \\[4pt]
    &\equiv (p \vee q) \vee \bigl[p \vee (\sim p \wedge \sim q)\bigr]
      && \text{Conmutativa} \\[4pt]
    &\equiv (p \vee q) \vee \bigl[(p \vee \sim p) \wedge (p \vee \sim q)\bigr]
      && \text{Distributiva} \\[4pt]
    &\equiv (p \vee q) \vee \bigl[\mathbf{V} \wedge (p \vee \sim q)\bigr]
      && \text{Complemento} \\[4pt]
    &\equiv (p \vee q) \vee (p \vee \sim q)
      && \text{Identidad} \\[4pt]
    &\equiv p \vee (q \vee \sim q)
      && \text{Asociativa + Conmutativa} \\[4pt]
    &\equiv p \vee \mathbf{V}
      && \text{Complemento} \\[4pt]
    &\equiv \mathbf{V}
      && \text{Identidad}
  \end{align*}
  \begin{tcolorbox}[colback=AzulClaro!50, colframe=AzulPizarra,
    arc=2pt, boxrule=0.7pt]
    \[ \boxed{\mathbf{V} \quad \text{(Tautolog\'{i}a)}} \]
  \end{tcolorbox}
\end{cajaejemplo}

\begin{cajaejemplo}[3 — Nivel admisión UNI: Simplificación con condicional]
  \textbf{Enunciado:} Determinar el equivalente más simple de:
  \[
    (p \vee q) \rightarrow (\sim p \wedge q)
  \]
  \noindent\textbf{Método:} Ley de la condicional, De Morgan, distributiva,
  complemento e identidad.
  \medskip

  \begin{align*}
    &(p \vee q) \rightarrow (\sim p \wedge q) \\[4pt]
    &\equiv \sim(p \vee q) \vee (\sim p \wedge q)
      && \text{Ley 7 (condicional)} \\[4pt]
    &\equiv (\sim p \wedge \sim q) \vee (\sim p \wedge q)
      && \text{De Morgan} \\[4pt]
    &\equiv \sim p \wedge (\sim q \vee q)
      && \text{Distributiva} \\[4pt]
    &\equiv \sim p \wedge \mathbf{V}
      && \text{Complemento} \\[4pt]
    &\equiv \sim p
      && \text{Identidad}
  \end{align*}
  \begin{tcolorbox}[colback=AzulClaro!50, colframe=AzulPizarra,
    arc=2pt, boxrule=0.7pt]
    \[ \boxed{\sim p} \]
  \end{tcolorbox}
\end{cajaejemplo}

% --- RESUMEN ---------------------------------------------------------------
\section{Resumen del tema}

\begin{cajaformula}
  \begin{center}
  \renewcommand{\arraystretch}{1.5}
  \begin{tabular}{@{}ll@{}}
    \toprule
    \rowcolor{AzulClaro!60}
    \textbf{Ley} & \textbf{Expresi\'{o}n principal} \\
    \midrule
    Doble negaci\'{o}n  & $\sim(\sim p) \equiv p$ \\
    Idempotencia        & $p \wedge p \equiv p$ \quad $p \vee p \equiv p$ \\
    Conmutativa         & $p \wedge q \equiv q \wedge p$ \\
    Asociativa          & $p \wedge(q \wedge r)\equiv(p \wedge q)\wedge r$ \\
    Distributiva        & $p \wedge(q \vee r)\equiv(p \wedge q)\vee(p \wedge r)$ \\
    De Morgan           & $\sim(p \wedge q)\equiv \sim p \vee \sim q$ \\
    Condicional         & $p \rightarrow q \equiv \sim p \vee q$ \\
    Bicondicional       & $p \leftrightarrow q \equiv (p \to q)\wedge(q \to p)$ \\
    Absorci\'{o}n       & $p \wedge(p \vee q)\equiv p$ \\
    Complemento         & $p \vee \sim p \equiv \mathbf{V}$ \quad $p \wedge \sim p \equiv \mathbf{F}$ \\
    Identidad           & $p \vee \mathbf{F}\equiv p$ \quad $p \wedge \mathbf{V}\equiv p$ \\
    \midrule
    Circuito serie      & $p \wedge q$ (ambos cierran) \\
    Circuito paralelo   & $p \vee q$ (al menos uno cierra) \\
    \bottomrule
  \end{tabular}
  \end{center}
\end{cajaformula}

\begin{cajatip}
  En el examen de admisión, la simplificación siempre termina en una
  proposición atómica, $\mathbf{V}$, $\mathbf{F}$, o una expresión de
  dos variables. Si tu resultado es más largo, revisa si aplicaste
  absorción o complemento correctamente.
\end{cajatip}

% FIN — teoria.tex
